% ============================================================================
% sec3_structural_model.tex — Cover-bidding model (full proofs in Appendix A)
% ============================================================================

\section{A Model of Cover Bidding in Procurement Auctions}
\label{sec:structural}

Cover bidding is well documented in prosecuted cartels
\citep{porter1999ohio, pesendorfer2000study, asker2010leniency,
marshall2012economics}, yet no existing model turns it into a source
of \emph{ex ante} testable predictions. We develop a model that does.
It distinguishes two cover-bidding regimes---complementary and
coordinated---and generates five empirical predictions that discipline
every subsequent test (Table~\ref{tab:predictions}). The model's
central implication is a detection impossibility result: under the
coordinated regime (Regime~2), cover bids cluster tightly above the
designated winner, so within-tender dispersion screens lose power
\emph{precisely when coordination is strongest}. A
participation-based screen sidesteps this vulnerability entirely.

The full model, assumptions, proofs, and structural likelihood are in
Appendix~\ref{app:proofs} and~\ref{app:structural_id}. Here we
present the setup, key results, and the predictions taken to data.

\subsection{Setup and Key Results}

A cartel controls a designated winner and deploys $m \geq 0$ cover
bidders (FL firms), each bidding above the cartel's winning bid
$b^*$. The cartel maximizes profit net of cover-bidding costs and
detection penalties, subject to a minimum-bidder constraint under
convite ($\underline{n} = 3$) and no constraint under preg\~ao.
Under standard assumptions on diminishing returns and increasing
detection risk (Appendix~\ref{app:proofs}), the optimal number of
cover bidders $m^*$ satisfies:
\begin{enumerate}
  \item[(i)] $m^*$ decreases in detection probability $\theta_k$;
  \item[(ii)] $m^*$ decreases in per-bidder cost $c_1$;
  \item[(iii)] $m^*$ is weakly larger under convite (corner solution
    from the minimum-bidder rule). However, the constraint-driven
    component dilutes the FL--price signal by mixing mandatory and
    voluntary deployments (Section~\ref{sec:mechanisms}).
\end{enumerate}

Two cover-bidding regimes are possible. Under \textbf{Regime~1}
(complementary), cover bids are uniformly distributed above $b^*$,
producing higher dispersion than genuine bids. Under \textbf{Regime~2}
(coordinated), cover bids follow a truncated normal centered near
$b^*$, producing \emph{lower} dispersion---a pattern also observed in
coordinated settings. BIC selects between
regimes using the mixture likelihood (Appendix~\ref{app:structural_id}).

The \textbf{market-selection prediction} follows from the cross-derivative
of cartel surplus: cover bidding is most profitable in competitive
markets (low HHI), where the marginal return to manufactured
competition is highest.

\subsection{Testable Predictions}
\label{sec:predictions}

Table~\ref{tab:predictions} maps the model's predictions to the
empirical tests reported below.

\begin{table}[htbp]
\centering
\caption{Model predictions and empirical tests}
\label{tab:predictions}
\small
\begin{tabular}{llll}
\toprule
\# & Prediction & Test & Section \\
\midrule
P1\label{pred:price_v6} & FL $\Rightarrow$ higher prices & OLS with FE & \ref{sec:reduced_form} \\
P2\label{pred:selection_v6} & FL adds to, not displaces, genuine bidders & Non-FL count & \ref{sec:reduced_form} \\
P3\label{pred:regime_v6} & Regime~2: $\sigma_c < \sigma_g$ & BIC model selection & \ref{sec:structural_estimation} \\
P4\label{pred:market_v6} & FL effect $\downarrow$ in HHI & Network split & \ref{sec:empirical_network} \\
P5\label{pred:exchange_v6} & FL residuals non-exchangeable & Bajari--Ye KS + pairwise & \ref{sec:bajari_ye} \\
\bottomrule
\end{tabular}
\end{table}

\noindent The model generates the predictions tested in
Sections~\ref{sec:results}--\ref{sec:mechanisms}; welfare
implications are in Appendix~\ref{app:extensions}.
